% Einheit 1 von 5 – Bruch als Anteil (bank/brueche-dezimalzahlen/e1.jsonl, zusammenbau v0.1)
%% TODO Zweigzeile Teil 1: was hier gelernt wird (Satz in Schülersprache aus dem Einheitstitel) – Einheit 1
\einheitenkopf[e1]{Einheit 1 von 5 · Bruch als Anteil}
\zweigzeile{ab Kl. 5 · P10 oft}
\setcounter{aufgabe}{9}

%% TODO \verfahren{…}: Verfahrensüberschrift in Sachsprache fehlt in der Bank (2.3 b)
%% TODO Titel: Ich-kann-Satz fehlt in der Bank (Platzhalter: Kettenname) – Nr. 10
%% TODO Anweisung: ein Satz über der ersten Teilaufgabe fehlt in der Bank – Nr. 10
\begin{aufgabe}{Anteil bestimmen}
%% TODO \rechenplatz{n}: Schrittzahl des Grundfalls fehlt in der Bank (nur bei fester Schreibform, 2.3 b)
\begin{teile}
\teil Sind alle Teile gleich groß? \janein

\bruchkreis{0}{6}
\teil Welcher Anteil des Streifens ist gefärbt?

\bruchrechteck{3}{5}
\leerfeld
\teil Färbe $\frac{5}{6}$ des Streifens.

\bruchrechteck{0}{6}
\teil Welcher Anteil vom Kreis sind die Sektoren A, B und D zusammen?

\kreisdiagramm{A/90,B/90,C/90,D/45,E/45}
\leerfeld
\teil Bei welcher Figur ist $\frac{2}{5}$ gefärbt? Kreuze an.\\ \kreuz{P}\\ \kreuz{Q}\\ \kreuz{R}

P \bruchrechteck{2}{7} \quad Q \bruchkreis{2}{5} \quad R \bruchrechteck{3}{5}
\teil In einer Gruppe sind $9$ Kinder, $4$ davon tragen eine Brille. Welcher Anteil trägt eine Brille? \leerfeld
\teil Markiert sind die Sektoren A und E. Welcher Anteil des Kreises ist markiert? \hfill (P10 2019 OS)\\ \kreuz{$\frac{2}{5}$}\\ \kreuz{$\frac{3}{9}$}\\ \kreuz{$\frac{2}{9}$}\\ \kreuz{$\frac{3}{5}$}

\kreisdiagramm{A/80,B/80,C/80,D/80,E/40}
\end{teile}
\end{aufgabe}

%% TODO \verfahren{…}: Verfahrensüberschrift in Sachsprache fehlt in der Bank (2.3 b)
%% TODO Titel: Ich-kann-Satz fehlt in der Bank (Platzhalter: Kettenname) – Nr. 11
%% TODO Anweisung: ein Satz über der ersten Teilaufgabe fehlt in der Bank – Nr. 11
\begin{aufgabe}{Bruchteil einer Größe}
%% TODO \rechenplatz{n}: Schrittzahl des Grundfalls fehlt in der Bank (nur bei fester Schreibform, 2.3 b)
\begin{teile}
\teil Eine Tüte enthält $24$ Bonbons, $\frac{1}{3}$ davon sind rot. Was ist das Ganze?
\teil $\frac{1}{4}$ von $32$
\teil $\frac{3}{4}$ von $36$
\teil $\frac{3}{5}$ von $45$ € \leerfeld[€]
\teil Ein Kuchen ist in $12$ Stücke geteilt, $\frac{2}{3}$ davon sind gegessen. Wie viele Stücke sind noch da? \leerfeld Stücke
\teil $\frac{2}{5}$ von $1{,}5$ kg – wie viel Gramm? \leerfeld[g]
\teil Ein Aquarium fasst $120$ l und ist zu $\frac{3}{4}$ gefüllt. Wie viele Liter fehlen bis zum Rand? \hfill (P10 2015 OS) \\ \leerfeld[l]
\end{teile}
\end{aufgabe}

%% TODO Titel: Ich-kann-Satz fehlt in der Bank (Platzhalter: Kettenname) – Nr. 12
%% TODO Anweisung: ein Satz über der ersten Teilaufgabe fehlt in der Bank – Nr. 12
\begin{aufgabe}{Ganzes aus Bruchteil}
\begin{teile}
\teil Ein Drittel der Kekse sind $8$ Stück. Wie viele Kekse sind es insgesamt? \leerfeld Kekse
\end{teile}
\end{aufgabe}

%% TODO Titel: Ich-kann-Satz fehlt in der Bank (Platzhalter: Kettenname) – Nr. 13
%% TODO Anweisung: ein Satz über der ersten Teilaufgabe fehlt in der Bank – Nr. 13
\begin{aufgabe}{Anteil bestimmen – Fehler finden, Begründen, Darstellungswechsel, Anwendung}
\begin{teile}
\teil Ben soll $\frac{2}{3}$ der Figur färben. Er färbt $2$ Kästchen. Finde den Fehler.

\bruchrechteck{0}{12}
\teil Eine Pizza wird einmal unter $4$ Kinder und einmal unter $6$ Kinder gerecht verteilt. Wann ist ein Stück größer? Begründe.
\teil Schreibe als Bruch: fünf Neuntel. \leerfeld
\teil Ein Fußballspiel dauert $90$ Minuten. Luca war $\frac{2}{3}$ der Spielzeit auf dem Platz. Wie viele Minuten saß er auf der Bank? \leerfeld[min]
\end{teile}
\end{aufgabe}
